% META: {"id": "1SPE-SECDEG-EX-006", "chapitre": "1SPE-SECOND-DEGRE", "type_objet": "exercice", "capacites": ["1SPE-SECOND-DEGRE-C2"], "capacites_codes": ["C2"], "methodes": ["M2"], "parcours": 1, "competences": ["calculer", "communiquer"], "duree_min": 12, "mode_creation": "ex_nihilo", "sources_inspiration": [], "parametres_sympy": {}, "fichier_tex": "chapitres/1SPE-SECOND-DEGRE/exercices/1SPE-SECDEG-EX-006.tex", "corrige_tex": "chapitres/1SPE-SECOND-DEGRE/corriges/1SPE-SECDEG-CO-006.tex", "status": "generated"}
% BEGIN-VERIFY
% from sympy import symbols, Rational, simplify
% x = symbols('x')
% # g(x) = -3x^2 + 12x - 5
% g = -3*x**2 + 12*x - 5
% alpha = Rational(-12, 2*(-3))   # -b/(2a) = 2
% beta = g.subs(x, alpha)
% assert alpha == 2
% assert beta == 7
% # forme canonique : -3*(x-2)^2 + 7
% assert simplify(g - (-3*(x-2)**2 + 7)) == 0
% # maximum : 7 en x = 2
% # tableau de variations : croissante sur (-inf, 2], decroissante sur [2, +inf)
% # f(0) = -5 (ordonnee a l'origine)
% assert g.subs(x, 0) == -5
% END-VERIFY
\begin{exercice}{1SPE-SECDEG-EX-006}{1}{12}
On considère la fonction $g$ définie sur $\mathbb{R}$ par $g(x) = -3x^2 + 12x - 5$.

\begin{enumerate}
  \item Calculer les coordonnées du sommet $S$ de la parabole représentative de $g$.
  \item Écrire $g$ sous forme canonique.
  \item Dresser le tableau de variations de $g$ sur $\mathbb{R}$. Préciser la valeur maximale de $g$ et l'abscisse en laquelle elle est atteinte.
  \item Calculer $g(0)$. Quelle est la signification géométrique de ce résultat ?
  \item Sans calcul, donner le signe de $g(4)$ en utilisant la symétrie de la parabole.
\end{enumerate}
\end{exercice}
